%% toc: P1 2 2.1
\section{Proofs for the separation instances of relative complement and intersection}
\label{sec:annex-p01-w02b-p01-004-separation-instances}

Throughout, the variables range over sets and the logic is classical
first-order logic with identity. ``By extensionality'' refers to clause 3 of
\cref{def:p01-w02-p01-005-membership}: two sets are equal iff they have the
same members; the direction from equality to sameness of members is the
substitutivity of identity. The bound-variable convention of
\cref{def:p01-w02b-p01-004-variables} is in force. The ordered-pair identity
of \cref{sec:set-prerequisites-02} states that
\(\langle a,b\rangle=\langle c,d\rangle\) iff \(a=c\) and \(b=d\).

\begin{annexitem}{annex:p01-w02b-p01-004-defining-formulas}{thm:p01-w02b-p01-004-defining-formulas}{The formulas define their intended sets}
\begin{proof}
\emph{Clause 1.} By the pairing axiom and extensionality, \(\{a,b\}\) is the
unique set \(c\) with \(\forall t\,(t\in c\leftrightarrow(t=a\lor t=b))\)
(\cref{sec:p01-w02-p01-005-operations-defs}). Assume
\(\theta_{\mathrm{pair}}(w,a,b)\). For an arbitrary set \(t\), the assumption
gives \(t\in w\leftrightarrow(t=a\lor t=b)\) and the pairing condition gives
\((t=a\lor t=b)\leftrightarrow t\in\{a,b\}\); hence
\(t\in w\leftrightarrow t\in\{a,b\}\), and \(w=\{a,b\}\) by extensionality.
Conversely, assume \(w=\{a,b\}\). For an arbitrary set \(t\), substitutivity
of identity gives \(t\in w\leftrightarrow t\in\{a,b\}\), and the pairing
condition turns the right side into \(t=a\lor t=b\); so
\(\theta_{\mathrm{pair}}(w,a,b)\). The second equivalence is the case \(b=a\)
of the first, because \(\{a\}\) is by definition \(\{a,a\}\).

\emph{Clause 2.} By definition \(\langle a,b\rangle=\{\{a\},\{a,b\}\}\), so by
the pairing condition, for every set \(s\),
\(s\in\langle a,b\rangle\leftrightarrow(s=\{a\}\lor s=\{a,b\})\). By clause 1,
\(s=\{a\}\leftrightarrow\theta_{\mathrm{pair}}(s,a,a)\) and
\(s=\{a,b\}\leftrightarrow\theta_{\mathrm{pair}}(s,a,b)\). Hence
\begin{equation}
  s\in\langle a,b\rangle\leftrightarrow
  \bigl(\theta_{\mathrm{pair}}(s,a,a)\lor\theta_{\mathrm{pair}}(s,a,b)\bigr)
  \qquad\text{for every set }s.
  \tag{\(\ast\)}
\end{equation}
Assume \(\theta_{\mathrm{op}}(z,a,b)\). For an arbitrary \(s\), the
assumption and \((\ast)\) give \(s\in z\leftrightarrow s\in\langle a,b\rangle\),
so \(z=\langle a,b\rangle\) by extensionality. Conversely, assume
\(z=\langle a,b\rangle\). For an arbitrary \(s\), substitutivity gives
\(s\in z\leftrightarrow s\in\langle a,b\rangle\), and \((\ast)\) turns this into
the matrix of \(\theta_{\mathrm{op}}(z,a,b)\); so
\(\theta_{\mathrm{op}}(z,a,b)\).

\emph{Clause 3.} By \cref{thm:p01-w02-p01-005-membership}(2),
\(v\in X\cap Y\leftrightarrow(v\in X\land v\in Y)\) for every set \(v\). If
\(\theta_{\cap}(u,X,Y)\), then \(v\in u\leftrightarrow v\in X\cap Y\) for every
\(v\), and \(u=X\cap Y\) by extensionality. If \(u=X\cap Y\), then
\(v\in u\leftrightarrow v\in X\cap Y\leftrightarrow(v\in X\land v\in Y)\) for
every \(v\) by substitutivity, which is \(\theta_{\cap}(u,X,Y)\).

\emph{Clause 4, left to right.} Assume \(\gamma_{\cap}(w,Q)\). Then there are
sets \(X,Y,p,u\) with \(X\in Q\), \(Y\in Q\), \(\theta_{\mathrm{op}}(p,X,Y)\),
\(\theta_{\cap}(u,X,Y)\) and \(\theta_{\mathrm{op}}(w,p,u)\). By clause 2,
\(p=\langle X,Y\rangle\) and \(w=\langle p,u\rangle\); by clause 3,
\(u=X\cap Y\). Substituting the first and third identities into the second gives
\(w=\langle\langle X,Y\rangle,X\cap Y\rangle\). So the sets \(X\) and \(Y\)
witness the right side.

\emph{Clause 4, right to left.} Assume \(X\in Q\), \(Y\in Q\) and
\(w=\langle\langle X,Y\rangle,X\cap Y\rangle\). The set
\(p\defeq\langle X,Y\rangle\) exists: pairing gives \(\{X\}=\{X,X\}\) and
\(\{X,Y\}\), and pairing again gives \(\{\{X\},\{X,Y\}\}\). The set
\(u\defeq X\cap Y\) exists by \cref{def:p01-w02-p01-005-operations}(2) and
\cref{thm:p01-w02-p01-005-separated-set}. By clause 2,
\(\theta_{\mathrm{op}}(p,X,Y)\); by clause 3, \(\theta_{\cap}(u,X,Y)\). By
substitutivity \(w=\langle p,u\rangle\), so \(\theta_{\mathrm{op}}(w,p,u)\) by
clause 2. Hence \(X,Y,p,u\) witness the four existential quantifiers of
\(\gamma_{\cap}(w,Q)\), and \(\gamma_{\cap}(w,Q)\) holds.
\end{proof}
\end{annexitem}

\begin{annexitem}{annex:p01-w02b-p01-004-relative-complement}{thm:p01-w02b-p01-004-relative-complement}{The relative complement as a separated set}
\begin{proof}
\emph{The formula and its side conditions.} Let \(\varphi(x,B)\) be the formula
\(\neg\,x\in B\), written \(x\notin B\) by
\cref{def:p01-w02-p01-005-membership}(1). It is the negation of an atomic
formula of \(\Lang_{\in}\), so it is a formula of \(\Lang_{\in}\) and contains
no defined notation. Its free variables are \(x\) and \(B\); the variable \(x\)
is the one being separated, and \(B\) is the single parameter, assigned the
given set \(B\). Let \(d\) be a variable distinct from \(x\) and from \(B\), to
serve as the variable for the formed set. The formula \(\neg\,x\in B\) contains
only the variables \(x\) and \(B\), so \(d\) does not occur in it, and in
particular is not free in it. The bound \(A\) is a given set. Thus all side
conditions of \cref{def:p01-w02-p01-005-bounded-formation} hold.

\emph{The separation instance.} For this formula the separation schema of
\cref{def:p01-w02-p01-005-existence-principles}(6) has the instance
\[
  \forall a\,\exists d\,\forall x\,
  \bigl(x\in d\leftrightarrow(x\in a\land x\notin B)\bigr),
\]
with \(B\) assigned the given set. By
\cref{thm:p01-w02-p01-005-separated-set}, applied with bound \(A\) and this
formula, there is exactly one set \(d_0\) such that
\[
  \forall x\,\bigl(x\in d_0\leftrightarrow(x\in A\land x\notin B)\bigr),
\]
and \(d_0\subseteq A\). By \cref{def:p01-w02-p01-005-operations}(3),
\(A\setminus B\) is the set \(\{x\in A:x\notin B\}\), which by
\cref{def:p01-w02-p01-005-bounded-formation} denotes this unique set \(d_0\).

\emph{The membership condition.} Let \(z\) be an arbitrary set. Instantiating
the universal quantifier at \(z\) gives
\(z\in A\setminus B\leftrightarrow(z\in A\land\varphi(z,B))\), where
\(\varphi(z,B)\) is \(\neg\,z\in B\): the formula \(\neg\,x\in B\) has no
quantifier, so replacing its one occurrence of \(x\) by \(z\) captures nothing.
Hence \(z\in A\setminus B\leftrightarrow(z\in A\land z\notin B)\), and since
\(z\) was arbitrary the displayed universal statement of the lemma holds.

\emph{Uniqueness and inclusion.} If \(D\) is a set with
\(\forall z\,(z\in D\leftrightarrow(z\in A\land z\notin B))\), then \(D\) has
the property that determines \(d_0\), and the uniqueness part of
\cref{thm:p01-w02-p01-005-separated-set} gives \(D=d_0=A\setminus B\). The
inclusion \(A\setminus B\subseteq A\) is \(d_0\subseteq A\).
\end{proof}
\end{annexitem}

\begin{annexitem}{annex:p01-w02b-p01-004-intersection-graph}{thm:p01-w02b-p01-004-intersection-graph}{Membership, single-valuedness and totality of the intersection graph}
\begin{proof}
Write \(T\defeq(Q\times Q)\times Q\). By
\cref{thm:p01-w02-p01-005-product-bound} with \(A=B=Q\), the members of
\(Q\times Q\) are exactly the sets \(\langle X,Y\rangle\) with \(X\in Q\) and
\(Y\in Q\); by the same lemma with \(A=Q\times Q\) and \(B=Q\), the members of
\(T\) are exactly the sets \(\langle m,c\rangle\) with \(m\in Q\times Q\) and
\(c\in Q\).

\emph{Clause 1, formation.} By \cref{def:p01-w02b-p01-004-graph} the side
conditions of bounded set formation hold for the bound \(T\) and the formula
\(\gamma_{\cap}(w,Q)\). By \cref{thm:p01-w02-p01-005-separated-set} there is
exactly one set \(G\) with
\begin{equation}
  \forall w\,\bigl(w\in G\leftrightarrow(w\in T\land\gamma_{\cap}(w,Q))\bigr),
  \tag{\(\dagger\)}
\end{equation}
and \(G\subseteq T\). By the definition of binary relation in
\cref{sec:set-prerequisites-03}, a subset of \((Q\times Q)\times Q\) is a binary
relation from \(Q\times Q\) to \(Q\).

\emph{Clause 1, membership, left to right.} Let \(w\in G\). By \((\dagger)\),
\(\gamma_{\cap}(w,Q)\) and \(w\in T\). By
\cref{thm:p01-w02b-p01-004-defining-formulas}(4) there are \(X,Y\in Q\) with
\(w=\langle\langle X,Y\rangle,X\cap Y\rangle\). Since \(w\in T\), there are
\(m\in Q\times Q\) and \(c\in Q\) with \(w=\langle m,c\rangle\). Then
\(\langle\langle X,Y\rangle,X\cap Y\rangle=\langle m,c\rangle\), and the
ordered-pair identity gives \(X\cap Y=c\). Hence \(X\cap Y\in Q\), and \(X\),
\(Y\) witness the right side of the membership condition.

\emph{Clause 1, membership, right to left.} Let \(X\in Q\), \(Y\in Q\),
\(X\cap Y\in Q\) and \(w=\langle\langle X,Y\rangle,X\cap Y\rangle\). Since
\(X\in Q\) and \(Y\in Q\), \(\langle X,Y\rangle\in Q\times Q\). Since moreover
\(X\cap Y\in Q\), the pair \(\langle\langle X,Y\rangle,X\cap Y\rangle\) is a
member of \(T\), and so is \(w\) by substitutivity. By
\cref{thm:p01-w02b-p01-004-defining-formulas}(4), \(\gamma_{\cap}(w,Q)\). By
\((\dagger)\), \(w\in G\).

\emph{Clause 2.} Let \(\langle\langle X,Y\rangle,u\rangle\in G\). By clause 1
there are \(X',Y'\in Q\) with \(X'\cap Y'\in Q\) and
\(\langle\langle X,Y\rangle,u\rangle=
\langle\langle X',Y'\rangle,X'\cap Y'\rangle\). The ordered-pair identity gives
\(\langle X,Y\rangle=\langle X',Y'\rangle\) and \(u=X'\cap Y'\); applied again
to the first equation it gives \(X=X'\) and \(Y=Y'\). Then \(X\cap Y\) and
\(X'\cap Y'\) have the same members by
\cref{thm:p01-w02-p01-005-membership}(2) and substitutivity, so
\(X'\cap Y'=X\cap Y\) by extensionality, and \(u=X\cap Y\).

For the second statement, let \(\langle m,u\rangle\in G\) and
\(\langle m,u'\rangle\in G\). By clause 1 applied to \(\langle m,u\rangle\),
there are \(X,Y\in Q\) with
\(\langle m,u\rangle=\langle\langle X,Y\rangle,X\cap Y\rangle\), and the
ordered-pair identity gives \(m=\langle X,Y\rangle\) and \(u=X\cap Y\). By
substitutivity \(\langle\langle X,Y\rangle,u'\rangle=\langle m,u'\rangle\in G\),
so the first statement gives \(u'=X\cap Y\). Hence \(u=u'\).

\emph{Clause 3, closure implies function.} Assume \(X\cap Y\in Q\) for all
\(X,Y\in Q\). By clause 1, \(G\subseteq(Q\times Q)\times Q\). Let
\(m\in Q\times Q\). Then \(m=\langle X,Y\rangle\) for some \(X,Y\in Q\). Put
\(b\defeq X\cap Y\); then \(b\in Q\) by the assumption. By clause 1, with the
witnesses \(X,Y\) and the hypothesis \(X\cap Y\in Q\), the set
\(\langle\langle X,Y\rangle,X\cap Y\rangle\) is a member of \(G\), and it equals
\(\langle m,b\rangle\) by substitutivity. So some \(b\in Q\) satisfies
\(\langle m,b\rangle\in G\). If \(b'\in Q\) also satisfies
\(\langle m,b'\rangle\in G\), then \(b'=b\) by clause 2. So for every
\(m\in Q\times Q\) there is exactly one \(b\in Q\) with
\(\langle m,b\rangle\in G\), and \(G\) is a function from \(Q\times Q\) to \(Q\)
by the definition of function in \cref{sec:set-prerequisites-03}. Its value at
\(\langle X,Y\rangle\) is the unique such \(b\), which is \(X\cap Y\).

\emph{Clause 3, function implies closure.} Assume that \(G\) is a function from
\(Q\times Q\) to \(Q\), and let \(X,Y\in Q\). Then
\(\langle X,Y\rangle\in Q\times Q\), so there is \(b\in Q\) with
\(\langle\langle X,Y\rangle,b\rangle\in G\). By clause 2, \(b=X\cap Y\). Hence
\(X\cap Y\in Q\) by substitutivity.
\end{proof}
\end{annexitem}

\begin{annexitem}{annex:p01-w02b-p01-004-intersection-on-subsets}{thm:p01-w02b-p01-004-intersection-on-subsets}{Intersection is an operation on the subsets of a set}
\begin{proof}
\emph{Closure.} The set \(Q=\mathcal P(S)\) exists by the power-set axiom, and
\(X\in Q\leftrightarrow X\subseteq S\) for every set \(X\). Let \(X,Y\in Q\), so
\(X\subseteq S\). For an arbitrary \(z\in X\cap Y\),
\cref{thm:p01-w02-p01-005-membership}(2) gives \(z\in X\), and \(X\subseteq S\)
gives \(z\in S\). So \(X\cap Y\subseteq S\), that is, \(X\cap Y\in Q\).

\emph{The function.} By the closure just proved and
\cref{thm:p01-w02b-p01-004-intersection-graph}(3), \(G^{Q}_{\cap}\) is a
function from \(Q\times Q\) to \(Q\) whose value at \(\langle X,Y\rangle\) is
\(X\cap Y\). By \cref{thm:p01-w02b-p01-004-intersection-graph}(1), a set \(w\)
is a member of \(G^{Q}_{\cap}\) iff there are \(X,Y\in Q\) with \(X\cap Y\in Q\)
and \(w=\langle\langle X,Y\rangle,X\cap Y\rangle\); the conjunct
\(X\cap Y\in Q\) follows from \(X,Y\in Q\) by closure, and \(X\in Q\) means
\(X\subseteq S\). So the members are exactly the triples
\(\langle\langle X,Y\rangle,X\cap Y\rangle\) with \(X,Y\subseteq S\).

\emph{Identification with \(G_{\cap}\).} The intersection operation \(G_{\cap}\)
of \cref{thm:p01-w02-p01-005-subset-algebra} is the set of members \(w\) of
\((Q\times Q)\times Q\) satisfying
\(\exists X\,\exists Y\,(X\in Q\land Y\in Q\land
w=\langle\langle X,Y\rangle,X\cap Y\rangle)\). By
\cref{thm:p01-w02b-p01-004-defining-formulas}(4) this condition is equivalent,
for every \(w\), to the \(\Lang_{\in}\)-formula \(\gamma_{\cap}(w,Q)\). So a set
is a member of \(G_{\cap}\) iff it is a member of \((Q\times Q)\times Q\)
satisfying \(\gamma_{\cap}(w,Q)\), iff it is a member of \(G^{Q}_{\cap}\); and
\(G_{\cap}=G^{Q}_{\cap}\) by extensionality.
\end{proof}
\end{annexitem}
